\documentclass[10pt,a4paper]{amsart}
\usepackage[margin=19mm]{geometry}
\usepackage{amsmath,amssymb,mathtools}
\usepackage[scr=rsfs,cal=euler]{mathalfa}
\usepackage{xfrac}
\usepackage[unicode,hidelinks,bookmarks=false]{hyperref}
\newcommand{\ie}{i.e.\ }
\newcommand{\cf}{cf.\ }
\newcommand{\resp}{respectively\ }
\newcommand{\Subsection}{\subsection}
\providecommand{\nobreakdash}{\nobreak\hbox{-}}
\setlength{\emergencystretch}{2em}
\multlinegap0pt
\usepackage{bm}
\usepackage{bbm}
\usepackage{dsfont}
\usepackage{xspace}
\usepackage{cancel}
\usepackage{mathtools}
\usepackage{amsthm}
\makeatletter
\@ifundefined{thm}{\newtheorem{thm}{Thm}}{}
\@ifundefined{lemma}{\newtheorem{lemma}[thm]{Lemma}}{}
\@ifundefined{proposition}{\newtheorem{proposition}[thm]{Proposition}}{}
\makeatother
\newcommand{\C}{\mathbb{C}}
\newcommand{\N}{\mathbb N}
\newcommand{\U}{\mathcal{U}}
\newcommand{\Z}{\mathbb Z}
\newcommand{\fg}{\mathfrak g}
\newcommand{\rD}{\mathrm{D}}
\title[Whittaker modules over the loop Virasoro algebra]{Whittaker modules over the loop Virasoro algebra}
\author{Zhiqiang Li, Shaobin Tan, Qing Wang}
\date{Source: arXiv:2509.24574v1 (CC BY 4.0)}
\begin{document}
\maketitle

\noindent\emph{Source and licence.} Whittaker modules over the loop Virasoro algebra. Version arXiv:2509.24574v1 (CC BY 4.0), distributed under the Creative Commons Attribution 4.0 International licence, as recorded in the arXiv licence metadata for this exact version and shown on the abstract page https://arxiv.org/abs/2509.24574v1. Full rights notice: licenses/arxiv-2509.24574-CC-BY-4.0.txt.\par\smallskip

\noindent\textit{Editorial scope.} Counted unit: the Proposition whose source label is \texttt{EXP-2N-1} in the article (source0.tex lines 600--602, source label \texttt{EXP-2N-1}), delivered together with the article's own complete original proof of it (source0.tex lines 603--710, one \texttt{proof} environment of 108 lines). No mathematics is written, repaired or re-derived: statement and proof are the article's own text line for line. Required context retained uncounted: lem:charCE (lines 359--363); lem:non-CE (lines 366--369); some-property-Whit (lines 516--520); EXP-2N (lines 552--554). Every definition, display and label of those lines is retained unchanged. Omitted from this package and not redistributed: 1--358, 364--365, 370--515, 521--551, 555--599, 711--1477. Nothing inside a retained excerpt is omitted or altered: every retained body is delivered exactly as the article's lines are written, so no omission receipt of any kind accompanies this package. Citations: the retained excerpts carry no citation command at all, so no bibliography entry of the article is transcribed and no reference is redistributed. Reference adaptation: none was needed and none was made; every reference inside the delivered unit resolves inside it, so no unresolved reference or citation is produced. Printed slips kept literal: no printed word, symbol, hypothesis or number of the counted statement or of its proof was altered. Layout and class: the article's own class is \texttt{amsart} with packages including amsmath, geometry, amsthm, graphics, tabularx, amssymb, shapepar, amscd, graphicx, xypic, hyperref, color, soul, bm; the wrapper is the lane's pinned \texttt{amsart} editorial wrapper at 10pt on A4, which restates the theorem-like environments the retained text needs and adds the article's own macro definitions that the retained text needs (\texttt{\textbackslash C}, \texttt{\textbackslash N}, \texttt{\textbackslash U}, \texttt{\textbackslash Z}, \texttt{\textbackslash fg}, \texttt{\textbackslash rD}), with extra wrapper packages (bm, bbm, dsfont, xspace, cancel, mathtools). The delivered numbering is the unit's own. Assets: no retained span contains a figure environment, a graphics-inclusion command, a PDF-page inclusion, a table or a TikZ picture; no image file is copied into this package, the package declares no asset dependency and no image file is redistributed with it. Licence: version arXiv:2509.24574v1 of this article is distributed under the Creative Commons Attribution 4.0 International licence, as recorded in the arXiv licence metadata for this exact version and shown on the abstract page https://arxiv.org/abs/2509.24574v1. A full rights notice is delivered at licenses/arxiv-2509.24574-CC-BY-4.0.txt. Characters: every retained excerpt is pure ASCII, so the delivered engine, XeLaTeX with its default Latin Modern text font, reports no missing character.
\begin{lemma}\label{lem:charCE}
For any $\phi\in\mathcal{A}^*$,
the annihilator ideal $\mathrm{Ann}(\phi)\neq 0$ if and only if $\phi\in\mathcal{E}$. Moreover,
$\mathcal{E}$ is an $\mathcal{A}$-submodule of $\mathcal{A}^*$.
\end{lemma}

\begin{lemma}\label{lem:non-CE}
For any $\phi\in\mathcal{A}^*\setminus\mathcal{E}$ and $f(t)\in\mathcal{A}$, we have $f(t).\phi\in\mathcal{E}$ if and
only if $f(t)=0$. Thus for any $\phi\in\mathcal{A}^*\setminus\mathcal{E}$ and non-zero element $f(t)\in\mathcal{A}$, \[f(t).\phi\in\mathcal{A}^*\setminus\mathcal{E}.\]
\end{lemma}

\begin{proposition}\label{some-property-Whit}
Let $v=\sum_{i=1}^{p}a_{i}v_{i}\in W(\fg_{\geq N},\psi)_{\psi}\setminus \C v_{\psi}$,
where $a_{1},a_{2},\dots,a_{p}\in\C^{\times}$ and $v_{1},v_{2},\dots,v_{p}\in B(\fg_{\geq N},\psi)$ are distinct elements.
Then \[\mathrm{lth}_{-1}(v_{i})=0\quad \text{for all}\quad i=1,\dots,p.\]
\end{proposition}

\begin{proposition}\label{EXP-2N}
If $\psi_{2N}\notin\mathcal{E}$, then $W(\fg_{\geq N},\psi)_{\psi}=\C v_{\psi}$.
\end{proposition}

\begin{proposition}\label{EXP-2N-1}
If $\psi_{2N}\in\mathcal{E}$ and $\psi_{2N-1}\notin\mathcal{E}$, then $W(\fg_{\geq N},\psi)_{\psi}=\C v_{\psi}$.
\end{proposition}

\begin{proof}
Suppose to the contrary that $W(\fg_{\geq N},\psi)_{\psi}\neq\C v_{\psi}$.
Then from Proposition \ref{some-property-Whit}, we know that
there exists some \[v=\sum_{i=1}^{P}\alpha_{i}\nu_{i}\in W(\fg_{\geq N},\psi)_{\psi}\setminus \C v_{\psi},\]
where $\alpha_{1},\dots,\alpha_{P}\in\C^{\times}$ and $\nu_{1},\dots,\nu_{P}\in B(\fg_{\geq N},\psi)\setminus\{v_{\psi}\}$
are distinct elements such that $\mathrm{lth}_{-1}(\nu_{i})=0$ for $i=1,\dots,P$.

Set \[\ell=\min\,\bigcup_{i=1}^{P}\mathcal{D}_{\mathrm{set}}(\nu_{i})\quad \text{and}\quad
\iota=\max\,\{\mathrm{lth}_{\ell}(\nu_{i})\mid i=1,\dots,P\}.\]
For convenience, we rewrite $v$ as
\[v=\sum_{i=1}^{p}a_{i}v_{i}+\sum_{j=1}^{Q}\beta_{j}\mu_{j}+\sum_{m=1}^{Q'}\gamma_{m}u_{m},\]
where the parameters satisfy:
\begin{itemize}
\item $p\in\Z_{+}$, $Q,Q'\in\N$, $a_{1},\dots,a_{p},\beta_{1},\dots,\beta_{Q},\gamma_{1},\dots,\gamma_{Q'}\in\C^{\times}$;
\item $v_{1},\dots,v_{p}\in B(\fg_{\geq N},\psi)$ are distinct elements such that $v_{1}\succ v_{2}\succ\cdots\succ v_{p}$
and $\mathrm{lth}_{\ell}(v_{i})=\iota$ for $i=1,\dots,p$;
\item $\mu_{1},\dots,\mu_{Q}\in B(\fg_{\geq N},\psi)$ are distinct elements
such that $\mathrm{lth}_{\ell}(\mu_{j})=\iota-1$ for $j=1,\dots,Q$;
\item $u_{1},\dots,u_{Q'}\in B(\fg_{\geq N},\psi)$ are distinct elements
such that $\mathrm{lth}_{\ell}(u_{m})\leq\iota-2$ for $m=1,\dots,Q'$.
\end{itemize}
Set $v_{1}=x_{N-1}\cdots x_{\ell+2}\rD(\ell+1,m_{1})^{n_{1}}\cdots\rD(\ell+1,m_{r})^{n_{r}}
\rD(\ell,k_{1})^{l_{1}}\cdots\rD(\ell,k_{s})^{l_{s}}v_{\psi}$,
where $x_{N-1}\in\U(\fg_{N-1}),\dots,x_{\ell+2}\in\U(\fg_{\ell+2})$ are nonzero elements, $r\in\N$,
$m_{r}<\cdots<m_{1}$, $n_{1},\dots,n_{r}\in\Z_{+}$, $k_{s}<\cdots<k_{1}$,
and $l_{1},\dots,l_{s}\in\Z_{+}$ with $l_{1}+\cdots+l_{s}=\iota$.
Write \[x_{\ell+1}:=\rD(\ell+1,m_{1})^{n_{1}}\cdots\rD(\ell+1,m_{r})^{n_{r}}.\]

Without loss of generality, we assume $\sum_{j=1}^{Q}\beta_{j}\mu_{j}\neq0$.
We rewrite $\sum_{j=1}^{Q}\beta_{j}\mu_{j}$ as
\[\sum_{j=1}^{Q}\beta_{j}\mu_{j}=\sum_{n\in\Z}b_{n}x_{N-1}\cdots x_{\ell+1}\rD(\ell+1,n)\rD(\ell,k_{1})^{l_{1}}\cdots\rD(\ell,k_{s})^{l_{s}-1}v_{\psi}+\sum_{j=1}^{q}\lambda_{j}w_{j},\]
where the parameters satisfy:
\begin{itemize}
\item[(c1)] there are only finitely many $b_{n}\in\C$ ($n\in\Z$) are non-zero;
\item[(c2)] $q\in\N$, $w_{1},\dots,w_{q}\in B(\fg_{\geq N},\psi)$ are distinct elements such that
$\mathrm{lth}_{\ell}(w_{j})=\iota-1$ for $j=1,\dots,q$;
\item[(c3)] for each $j=1,\dots,q$, either $\mathcal{T}_{\ell}(w_{j})\neq \mathcal{T}_{\ell}(v_{+})$ or
$\mathcal{T}_{\ell}(w_{j})=\mathcal{T}_{\ell}(v_{+})$ but $w_{j}\not\in\{x_{N-1}\cdots x_{\ell+1}\rD(\ell+1,n)\rD(\ell,k_{1})^{l_{1}}\cdots\rD(\ell,k_{s})^{l_{s}-1}v_{\psi}\mid n\in\Z\}$,
where \[v_{+}=x_{N-1}\cdots x_{\ell+2}\rD(\ell+1,m_{1})^{n_{1}}\cdots\rD(\ell+1,m_{r})^{n_{r}}
\rD(\ell,k_{1})^{l_{1}}\cdots\rD(\ell,k_{s})^{l_{s}-1}v_{\psi}.\]
\end{itemize}

\vspace{2mm}

For $k\in\Z$, we have
$$\big(\rD(2N-1-\ell,k)-\psi_{2N-1-\ell}(t^{k})\big).v_{1}=l_{s}(2\ell-2N+1)
\psi_{2N-1}(t^{k+k_{s}})v_{+}+\sum_{i=1}^{p''}a_{i}''v_{i}'',$$
where $a_{1}'',\dots,a_{p''}''\in\C^{\times}$ and $v_{+},v_{1}'',\dots,v_{p''}''\in B(\fg_{\geq N},\psi)$
are distinct elements such that $v_{i}''\prec v_{+}$ for $i=1,\dots,p''$.


\vspace{2mm}
Similar to the proof of Proposition \ref{EXP-2N}, we know that there exists $1\leq p_{0}\leq p$ such that
\[v_{i}=x_{N-1}\cdots x_{\ell+2}\rD(\ell+1,m_{1})^{n_{1}}\cdots\rD(\ell+1,m_{r})^{n_{r}}
\rD(\ell,k_{1})^{l_{1}}\cdots\rD(\ell,k_{s})^{l_{s}-1}\rD(\ell,k_{s_{i}})v_{\psi}\]
for $i=2,\dots,p_{0}$, where $k_{s}>k_{s_{2}}>\cdots>k_{s_{p_{0}}}$ and that
\[v_{p_{0}+1}\prec x_{N-1}\cdots x_{\ell+2}\rD(\ell+1,m_{1})^{n_{1}}\cdots\rD(\ell+1,m_{r})^{n_{r}}
\rD(\ell,k_{1})^{l_{1}}\cdots\rD(\ell,k_{s})^{l_{s}-1}\rD(\ell,m)v_{\psi}\]
for any $m\in\Z$ with $m<k_{s}$. Then we obtain that, for $k\in\Z$,
$$\big(\rD(2N-1-\ell,k)-\psi_{2N-1-\ell}(t^{k})\big).\sum_{i=1}^{p}a_{i}v_{i}=b_{0}v_{+}+\sum_{i=1}^{p'}b_{i}'v_{i}',$$
where $b_{0}:=l_{s}a_{1}(2\ell-2N+1)\psi_{2N-1}(t^{k+k_{s}})+\sum_{i=2}^{p_{0}}a_{i}(2\ell-2N+1)\psi_{2N-1}(t^{k+k_{s_{i}}}),\
b_{1}',\dots,b_{p'}'\in\C$ and $v_{1}',\dots,v_{p'}'\in B(\fg_{\geq N},\psi)$ are distinct elements such that
$v_{i}'\prec v_{+}$ for $i=1,\dots,p'$.

\vspace{2mm}

\noindent {\bf Claim 1}.\quad For $k\in\Z$, we have
$$\big(\rD(2N-1-\ell,k)-\psi_{2N-1-\ell}(t^{k})\big).\sum_{j=1}^{Q}\beta_{j}\mu_{j}=c_{0}v_{+}+\sum_{j=1}^{q'}c_{j}'w_{j}',$$
where $c_{0}:=\sum_{n\in\Z}b_{n}(2\ell+2-2N)\big(1+\sum_{i=1}^{r}n_{i}\delta_{n,m_{i}}\big)\psi_{2N}(t^{k+n}),\ c_{1}',\dots,
c_{q'}'\in\C$
and $w_{1}',\dots,w_{q'}'\in B(\fg_{\geq N},\psi)\setminus\{v_{+}\}$.

Note that $\ell+1\geq1$, which implies that we only need to collect the elements
appearing in $\sum_{j=1}^{Q}\beta_{j}\mu_{j}$ that have the form
\[x_{N-1}\cdots x_{\ell+1}\rD(\ell+1,n)\rD(\ell,k_{1})^{l_{1}}\cdots\rD(\ell,k_{s})^{l_{s}-1}v_{\psi},\quad n\in\Z.\]
%We can deduce that $\rD(\ell+1,n)$ ($n\in\Z$) occurs only once.
Then from (c1)-(c3), we obtain Claim 1.




\noindent {\bf Claim 2}.\quad For $k\in\Z$, we have
$$\big(\rD(2N-1-\ell,k)-\psi_{2N-1-\ell}(t^{k})\big).\sum_{m=1}^{Q'}\gamma_{m}u_{m}=\sum_{j=1}^{Q''}\gamma_{j}'z_{j},$$
where $\gamma_{1}',\dots,\gamma_{Q''}'\in\C$ and $z_{1},\dots,z_{Q''}\in B(\fg_{\geq N},\psi)\setminus\{v_{+}\}$.

In fact, we have $\mathrm{lth}_{\ell}(z_{j})\leq\iota-2$ for all $j=1,\dots,Q''$.
Thus $z_{j}\neq v_{+}$ for all $j=1,\dots,Q''$.

\vspace{3mm}

Since $\big(\rD(2N-1-\ell,k)-\psi_{2N-1-\ell}(t^{k})\big).v=0$ for all $k\in\Z$, this forces
$b_{0}+c_{0}=0$ for all $k\in\Z$.
That is,
$$\begin{aligned}
\Big[\Big(\big(l_{s}a_{1}(2\ell-2N+1)t^{k_{s}}+\sum_{i=2}^{p_{0}}a_{i}(2\ell-2N+1)t^{k_{s_{i}}}\big).\psi_{2N-1}\Big)
+\\ \quad \Big(\big(\sum_{n\in\Z}b_{n}(2\ell+2-2N)\big(1+\sum_{i=1}^{r}n_{i}\delta_{n,m_{i}}\big)t^{n}\big).\psi_{2N}\Big)\Big](t^{k})=0
\end{aligned}$$
for all $k\in\Z$, which means
$$\begin{aligned}
\Big(l_{s}a_{1}(2\ell-2N+1)t^{k_{s}}+\sum_{i=2}^{p_{0}}a_{i}(2\ell-2N+1)t^{k_{s_{i}}}\Big).\psi_{2N-1}
+\\ \quad \Big(\sum_{n\in\Z}b_{n}(2\ell+2-2N)\big(1+\sum_{i=1}^{r}n_{i}\delta_{n,m_{i}}\big)t^{n}\Big).\psi_{2N}=0
\end{aligned}$$

Note that \[l_{s}a_{1}(2\ell-2N+1)t^{k_{s}}+\sum_{i=2}^{p_{0}}a_{i}(2\ell-2N+1)t^{k_{s_{i}}}\in\mathcal{A}\]
is non-zero element. Since $\psi_{2N}\in\mathcal{E}$ and $\psi_{2N-1}\notin\mathcal{E}$,
from Lemmas \ref{lem:charCE} and \ref{lem:non-CE}, we get a contradiction.
Note that even if all $b_{n}=0$ for $n\in\Z$, we also can get the contradiction. This completes the proof.
\end{proof}

\end{document}
